\section{Ventajas y cuellos de botella del Transformer: por qué la generalidad es costosa}

La sección anterior propuso una arquitectura unificada; esta sección explica por qué la investigación parte del Transformer. La attention del Transformer no impone conexiones locales, la posición suele tratarse como una característica de entrada y no como una topología de cómputo fija, y la mayor parte del cómputo consiste en multiplicaciones de matrices regulares. Estas propiedades hacen que se transfiera con más facilidad que la convolución a estructuras desconocidas; el costo es que, cuanto mayor es el raw input, más costosas resultan las comparaciones por pares de la self-attention.

\lecturefigure{slide-03-improving-transformers.jpg}{El inductive bias general del Transformer y el cuello de botella por el tamaño de la entrada.}{Video oficial actual de Stanford Online, 00:05:10--00:06:01.}

\begin{knowledgebox}{Glosario: las interfaces centrales de esta clase}
\begin{tabularx}{\textwidth}{L{0.23\textwidth}X}
\toprule
Término & Definición precisa \\
\midrule
non-local & En principio, cualquier elemento de la entrada puede atender directamente a cualquier otro, sin restringirse de antemano a una vecindad fija. \\
position as feature & Las coordenadas se codifican en las características del token, en lugar de limitar rígidamente las conexiones con kernels de convolución o aristas de un grafo. \\
latent array & Un número fijo o relativamente pequeño de vectores aprendibles que sirve como espacio de trabajo interno para procesar entradas grandes. \\
cross-attention & La query y las key/value provienen de arreglos distintos; aquí, el latent consulta la entrada original. \\
output query & Consulta del decoder que especifica «qué predecir, en qué posición, en qué modalidad o para qué tarea». \\
Fourier feature & Usa senos y cosenos de múltiples frecuencias para mapear coordenadas continuas a características de alta frecuencia aprendibles. \\
byte-level language & Modela directamente símbolos de bajo nivel, como los bytes UTF-8, sin depender de un tokenizer de subpalabras. \\
optical flow & Salida estructurada densa que predice, para cada píxel de fotogramas contiguos, un vector de movimiento bidimensional. \\
\bottomrule
\end{tabularx}
\end{knowledgebox}

La attention estándar proyecta la entrada $X\in\mathbb{R}^{M\times d}$ en query, key y value:

\[
Q=XW_Q,\qquad K=XW_K,\qquad V=XW_V,
\qquad
\operatorname{Attn}(X)=\operatorname{softmax}\!\left(\frac{QK^\top}{\sqrt{d_k}}\right)V.
\]

Aquí, $M$ es el número de elementos de la entrada; $d$ es la dimensión de la representación; $W_Q,W_K,W_V$ son las matrices de proyección; $d_k$ es la dimensión de las key. La matriz $QK^\top$ tiene forma $M\times M$, por lo que el cálculo y el almacenamiento de los puntajes de atención crecen con $M^2$.

\lecturefigure{slide-04-standard-qkv-attention.jpg}{Attention QKV estándar: cada elemento de la entrada se compara con todas las key.}{Video oficial actual de Stanford Online, 00:06:02--00:08:01.}

\begin{importantbox}{El costo de la self-attention según el tamaño de la entrada}
Si se ignoran las constantes, el costo central de los puntajes en una capa de attention es de aproximadamente $\mathcal{O}(M^2d)$, y la memoria de la matriz de atención, de aproximadamente $\mathcal{O}(M^2)$. Si se pasa de imágenes de $224^2$ píxeles a video o a flujos multimodales de alta resolución, el crecimiento de $M$ supera rápidamente al aumento de profundidad del modelo.
\end{importantbox}

\subsection{Non-locality: una forma de conexión costosa pero general}

La convolución presupone que los píxeles vecinos están más relacionados y puede procesar cuadrículas de manera eficiente con kernels locales. La attention, en cambio, permite que elementos distantes interactúen directamente, por lo que puede manejar conjuntos, correspondencias a larga distancia y topologías desconocidas. El ponente considera la non-locality un elemento esencial de la generalidad, no un desperdicio que haya que eliminar por completo.

\lecturefigure{slide-05-why-non-locality.jpg}{La non-local attention no fija de antemano una vecindad local; la convolución, en cambio, limita las conexiones a la adyacencia de la cuadrícula.}{Video oficial actual de Stanford Online, 00:08:02--00:08:37.}

\teachervoice{El juicio del docente no es que «la attention siempre supera a la convolution», sino que ambas ocupan puntos de Pareto distintos: cuando se conoce la estructura de la imagen, la convolución suele ser más rápida y requerir menos datos; ante cámaras de eventos, células, nubes de puntos o arreglos científicos de estructura desconocida, es posible que ni siquiera sepamos cuál es la vecindad de convolución correcta.}

\subsection{Position as feature: dejar las coordenadas al aprendizaje en lugar de fijarlas en las conexiones}

Esta sección se ocupa de la condición que acompaña a la non-locality: si la attention es por naturaleza insensible al orden de la entrada, el modelo necesita recibir la posición para reconstruir relaciones espaciales, temporales o secuenciales. Perceiver concatena la codificación de coordenadas con el contenido, de modo que la arquitectura sabe «dónde está esto», pero sin imponer «con quién puede conectarse». Esto es más útil que prescindir por completo de la posición y, a la vez, supone menos que un kernel local fijo.

\lecturefigure{slide-06-why-featurize-position.jpg}{Codificar la posición como característica permite que la attention aprenda relaciones espaciales en lugar de usar reglas de conexión rígidas.}{Video oficial actual de Stanford Online, 00:08:38--00:09:29.}

\begin{warningbox}{Un sesgo débil no equivale a independencia de la posición}
Si la tarea exige una salida espacial, el modelo todavía debe distinguir coordenadas. Sin posición, distintos ordenamientos se vuelven indistinguibles; dar la posición y permitir attention global es «debilitar la estructura local», no «eliminar la estructura».
\end{warningbox}

\subsection{La tensión entre escalabilidad y generality}

La clase sitúa la convnet y el Transformer en un mismo diagrama conceptual: la convnet está muy optimizada para cuadrículas y su complejidad puede ser aproximadamente lineal, pero admite un rango de entradas más estrecho; el Transformer se adapta a conjuntos generales, pero carga con una attention cuadrática. El objetivo de diseño de Perceiver es conservar la no localidad, la codificación de la posición como característica y el uso compartido de pesos, y al mismo tiempo reformular la complejidad respecto al tamaño de la entrada.

\lecturefigure{slide-07-scalability-vs-generality.jpg}{El compromiso entre ConvNet y Transformer en los supuestos sobre la estructura de la entrada y en la complejidad.}{Video oficial actual de Stanford Online, 00:09:30--00:09:41.}

\begin{knowledgebox}{Lectura de la figura: no es una «ruta evolutiva» de izquierda a derecha}
Los dos extremos de la figura representan objetivos de optimización distintos. Si la tarea se limita a clasificación de imágenes y los datos son escasos, el sesgo local de la convnet puede ser preferible; si la estructura de la entrada cambia con frecuencia, la interfaz unificada de Transformer/Perceiver puede reducir el costo de desarrollo. La evaluación debe considerar a la vez la calidad, la cantidad de datos, la velocidad y la reutilización entre dominios.
\end{knowledgebox}

\subsection{Resumen de la sección}

La generalidad del Transformer proviene de la non-local attention, de la codificación de la posición como característica y del cómputo matricial compartido, pero la self-attention sobre el raw input tiene un costo de $M^2$. Lo que Perceiver busca resolver es precisamente «conservar la interfaz y reducir el costoso espacio de trabajo».
